\section{2012级工科数学分析下 A卷}

\subsection{资料来源}
\begin{itemize}
  \item 试题：2012级计算机工科数学分析下期末考试题 A 卷，DOC 正文已抽取；公式对象已按 Word 公式图片逐项恢复。
  \item 公式图片审计：\texttt{tmp/past-exam-equations/2012-A/manifest.csv}。
\end{itemize}

\subsection{试题}
诚信应考，考试作弊将带来严重后果！

华南理工大学期末考试

《工科数学分析》2012--2013 第二学期期末考试试卷（A）

注意事项：
\begin{enumerate}
  \item 考前请将密封线内填写清楚；
  \item 所有答案请直接答在试卷上（或答题纸上）；
  \item 考试形式：闭卷；
  \item 本试卷共 5 个大题，满分 100 分，考试时间 120 分钟。
\end{enumerate}

\subsubsection*{一、单项选择题（每小题 3 分，共 15 分）}
\begin{enumerate}
  \item 下述函数中有可能成为二阶微分方程
  \[
  f(x,y,y'')=0
  \]
  的通解的是（\quad）。

  A. $C_1x+C_2y=0$；\quad B. $y-C_1^2=e^x+C_2$；\quad
  C. $C_1y=C_2e^{C_3x}$；\quad D. $y=C_1\ln x+C_2x+C_3$。

  \item 设 $f(x,y)$ 是可微函数，且
  \[
  f(0,0)=0,\qquad f_x(0,0)=a,\qquad f_y(0,0)=b,
  \]
  令 $\varphi(t)=f(t,f(t,t))$，则 $\varphi'(0)=$（\quad）。

  A. $a$；\quad B. $a+b(a+b)$；\quad C. $a+1$；\quad D. $\dfrac{a}{1-b}$。

  \item
  \[
  \int_0^1dy\int_{-y}^{\sqrt{2y-y^2}} f(x,y)\,dx
  \]
  交换积分顺序后，正确的结果是（\quad）。

  A. $\displaystyle \int_{-1}^0dx\int_{-x}^1f(x,y)\,dy+\int_0^1dx\int_{1-\sqrt{1-x^2}}^1f(x,y)\,dy$；

  B. $\displaystyle \int_0^1dx\int_{1-\sqrt{1-x^2}}^1f(x,y)\,dy$；

  C. $\displaystyle \int_{-1}^0dx\int_{-x}^1f(x,y)\,dy$；

  D. $\displaystyle \int_{-1}^1dx\int_{1-\sqrt{1-x^2}}^1f(x,y)\,dy$。

  \item 设 $C$ 是右半圆周
  \[
  x^2+y^2=a^2,\quad x\ge0,\quad a>0,
  \]
  则曲线积分
  \[
  \int_C(x+y)\,ds
  \]
  （\quad）。

  A. $0$；\quad B. $2a^2$；\quad C. $a^2$；\quad D. $4a$。

  \item 设 $f(x)=x^2\ (0\le x<1)$，而
  \[
  S(x)=\sum_{n=1}^{\infty}b_n\sin n\pi x\quad(-\infty<x<+\infty),
  \]
  其中
  \[
  b_n=2\int_0^1f(x)\sin n\pi x\,dx\quad(n=1,2,\cdots),
  \]
  则 $S\!\left(-\dfrac12\right)$（\quad）。

  A. $-\dfrac12$；\quad B. $-\dfrac14$；\quad C. $\dfrac14$；\quad D. $\dfrac12$。
\end{enumerate}

\subsubsection*{二、填空题（每小题 3 分，共 15 分）}
\begin{enumerate}
  \item 微分方程
  \[
  yy''-y'^2=0
  \]
  的通解为 \underline{\hspace{4cm}}。
  \item 设
  \[
  u=x+(y-1)\arcsin\frac{x}{y},
  \]
  则 $\left.\dfrac{\partial u}{\partial x}\right|_{(1,2)}=$ \underline{\hspace{3cm}}。
  \item 设区域
  \[
  D:0\le y\le x,\quad 0\le x\le\pi,
  \]
  则二重积分
  \[
  \iint_D \sqrt{1-\sin^2x}\,dx\,dy=
  \]
  \underline{\hspace{3cm}}。
  \item 已知曲线 $C$ 为
  \[
  x^2+y^2=ax,
  \]
  则曲线积分
  \[
  \int_C\sqrt{x^2+y^2}\,ds=
  \]
  \underline{\hspace{3cm}}。
  \item 幂级数
  \[
  \sum_{n=1}^{\infty}\frac{(x-2)^n}{n4^n}
  \]
  的收敛域为 \underline{\hspace{4cm}}。
\end{enumerate}

\subsubsection*{三、计算题（每小题 10 分，共 50 分）}
\begin{enumerate}
  \item 计算三重积分
  \[
  I=\iiint_\Omega (x^2+y^2+z^2)\,dV,
  \]
  其中 $\Omega:x^2+y^2+z^2\le2z$。
  \item 设 $\Sigma$ 是锥面 $z=\sqrt{x^2+y^2}$ 被平面 $z=0$ 及 $z=1$ 所截部分的外侧，计算第二类曲面积分
  \[
  I=\iint_\Sigma xy\,dy\,dz+y\,dz\,dx+(z^2-2z)\,dx\,dy.
  \]
  \item 求幂级数
  \[
  \sum_{n=0}^{\infty}(2n+1)x^n
  \]
  的和函数，并据此求数项级数
  \[
  \sum_{n=0}^{\infty}\frac{2n+1}{2^n}
  \]
  的值。
  \item 求 $a,b$ 的值，使得包含圆周
  \[
  (x-1)^2+y^2=1
  \]
  在其内部的椭圆
  \[
  \frac{x^2}{a^2}+\frac{y^2}{b^2}=1\quad(a>0,\ b>0,\ a\ne b)
  \]
  有最小的面积。
  \item 求微分方程
  \[
  y''-y=2xe^x
  \]
  的通解。
\end{enumerate}

\subsubsection*{四、证明题（本题 10 分）}
证明数项级数
\[
\sum_{n=1}^{\infty}\sin(\pi\sqrt{n^2+1})
\]
条件收敛。

\subsubsection*{五、应用题（本题 10 分）}
设某山峰可由曲面
\[
z=5-x^2-2y^2
\]
表示。位于点
\[
\left(\frac12,-\frac12,\frac{17}{4}\right)
\]
处的登山者发现其供氧面具漏气，必须返回。他应该沿哪个方向才能最快到达山底？

\subsection{答案与解析}

\subsubsection*{一、单项选择题}
\begin{enumerate}
  \item 二阶微分方程通解通常含两个独立任意常数，只有 B 项含两个独立常数且能表示 $y$，故选 B。
  \item 链式法则给出
  \[
  \varphi'(0)=f_x(0,0)+f_y(0,0)\bigl(f_x(0,0)+f_y(0,0)\bigr)=a+b(a+b).
  \]
  故选 B。
  \item 区域由 $-y\le x$ 与 $x^2+(y-1)^2\le1$ 组成，故选 A。
  \item 右半圆可取 $x=a\cos t,\ y=a\sin t,\ -\dfrac\pi2\le t\le\dfrac\pi2$。于是
  \[
  \int_C(x+y)\,ds=a^2\int_{-\pi/2}^{\pi/2}(\cos t+\sin t)\,dt=2a^2.
  \]
  故选 B。
  \item 正弦级数对应奇延拓且周期为 $2$，故
  \[
  S\!\left(-\frac12\right)=-f\!\left(\frac12\right)=-\frac14.
  \]
  故选 B。
\end{enumerate}

\subsubsection*{二、填空题}
\begin{enumerate}
  \item 设 $p=y'$，由
  \[
  yy''-(y')^2=0
  \]
  得 $\left(\dfrac{y'}{y}\right)'=0$，故
  \[
  y=Ce^{kx}.
  \]
  \item
  \[
  u_x=(y-1)\frac1{\sqrt{1-(x/y)^2}}\cdot\frac1y
  =\frac{y-1}{\sqrt{y^2-x^2}}.
  \]
  在 $(1,2)$ 处为 $\dfrac1{\sqrt3}$。
  \item 因 $x\in[0,\pi]$，有 $\sqrt{1-\sin^2x}=|\cos x|$。所以
  \[
  \iint_D|\cos x|\,dx\,dy=\int_0^\pi x|\cos x|\,dx=\pi.
  \]
  \item 圆 $x^2+y^2=ax$ 上可用极坐标 $r=a\cos\theta$，结果为
  \[
  2a^2.
  \]
  \item 半径 $R=4$，端点 $x=-2,6$。在 $x=6$ 时为 $\sum 1/n$ 发散；在 $x=-2$ 时为 $\sum(-1)^n/n$ 收敛，故收敛域为
  \[
  [-2,6).
  \]
\end{enumerate}

\subsubsection*{三、计算题}
\begin{enumerate}
  \item 区域是球心 $(0,0,1)$、半径 $1$ 的球。令 $z=w+1$，则
  \[
  x^2+y^2+z^2=x^2+y^2+w^2+1+2w.
  \]
  对称性使 $2w$ 积分为 $0$，故
  \[
  I=\iiint_{r\le1}r^2\,dV+\iiint_{r\le1}1\,dV
  =\frac{4\pi}{5}+\frac{4\pi}{3}=\frac{32\pi}{15}.
  \]
  \item 外侧锥面 $z=r$ 对应远离 $z$ 轴且 $z$ 分量为负的法向
  \[
  (z_x,z_y,-1)=\left(\frac{x}{r},\frac{y}{r},-1\right).
  \]
  因而
  \[
  dy\,dz=\frac{x}{r}\,dx\,dy,\quad dz\,dx=\frac{y}{r}\,dx\,dy,\quad dx\,dy=-dx\,dy.
  \]
  在 $D:x^2+y^2\le1$ 上积分得
  \[
  I=\iint_D\left(\frac{x^2y}{r}+\frac{y^2}{r}-r^2+2r\right)\,dx\,dy
  =\frac{\pi}{3}-\frac{\pi}{2}+\frac{4\pi}{3}
  =\frac{7\pi}{6}.
  \]
  \item 当 $|x|<1$，
  \[
  \sum_{n=0}^{\infty}(2n+1)x^n
  =2\sum_{n=0}^{\infty}nx^n+\sum_{n=0}^{\infty}x^n
  =\frac{2x}{(1-x)^2}+\frac1{1-x}
  =\frac{1+x}{(1-x)^2}.
  \]
  令 $x=\dfrac12$，得
  \[
  \sum_{n=0}^{\infty}\frac{2n+1}{2^n}=6.
  \]
  \item 与 2014 A 卷同型，最小面积椭圆为
  \[
  a=2,\qquad b=\sqrt2.
  \]
  \item 齐次解为 $C_1e^x+C_2e^{-x}$。右端为一次多项式乘 $e^x$，且 $1$ 是一重特征根，设
  \[
  y_p=e^x(Ax^2+Bx).
  \]
  代入得 $A=\dfrac12,\ B=-\dfrac12$，故
  \[
  y=C_1e^x+C_2e^{-x}+\frac12e^x(x^2-x).
  \]
\end{enumerate}

\subsubsection*{四、证明题}
因
\[
\sqrt{n^2+1}=n+\frac1{2n}+O\!\left(\frac1{n^3}\right),
\]
故
\[
\sin(\pi\sqrt{n^2+1})
=(-1)^n\sin\left(\frac{\pi}{2n}+O\!\left(\frac1{n^3}\right)\right)
\sim (-1)^n\frac{\pi}{2n}.
\]
原级数由交错级数判别法收敛；而绝对值级数与 $\sum 1/n$ 同阶发散，故条件收敛。

\subsubsection*{五、应用题}
山峰高度函数为 $z=5-x^2-2y^2$。最快下山方向为负梯度方向：
\[
-\nabla z=(2x,4y).
\]
在 $\left(\dfrac12,-\dfrac12\right)$ 处，
\[
-\nabla z=(1,-2).
\]
故应沿水平面内方向 $(1,-2)$ 前进，单位方向为
\[
\frac{(1,-2)}{\sqrt5}.
\]
